% Appendices (drafted 2026-11-18, D185; R7-3; docs/paper_outline.md appendices A and B). Every claim cites its
% source in a comment; every number is \pv{...} or sits in a generated table.
\section{Reproducibility}
\label{app:repro}

% D185 (R7-3 A; paper/RELEASE.yaml, the pin recorded in the commit after the sources); D150 (every value, table and
% figure generated from committed files; paperlib.Sources refuses untracked or modified sources).
\paragraph{Code, data and the pinned version.}
The code, configurations, evaluation sets, completion logs and analysis outputs are public at \pv{rel.repo}.
This paper was built from commit \pv{rel.commit}.
Every value in the text is a generated macro, and every table and figure is generated, from files committed in that
repository; \path{paper/generated/SOURCES.md} lists each artifact's source files with their git blob hashes, and the
build refuses a source that git does not track or that has uncommitted changes.

% D150 (scripts/paper/build.py); tests/test_paper.py (the committed values and tables equal a fresh build) and
% tests/test_figures.py (the figure plots the committed values).
\paragraph{Regenerating the headline figure and the tables.}
One command, run on a CPU from the pinned commit, regenerates Figure~\ref{fig:headline}, every table (for example
Table~\ref{tab:cost}) and every value in the text from the saved result files:
\begin{verbatim}
uv run python scripts/paper/build.py
\end{verbatim}
It writes them to \path{paper/generated/}, and \texttt{tectonic main.tex}, run in \path{paper/}, then compiles the
paper.
The repository's tests check that the committed values and tables equal a fresh build.

% CLAUDE.md component 6 (tape and replay); D72, D77 (R2-1: the training-time fingerprint and the reload check); D78
% (pinned bases); .gitignore (adapter weights are local only).
\paragraph{Logs and checkpoints.}
The completion logs of the reported training runs and evaluations are in the repository.
Each line holds the problem's seed and generator version, the completion, the verifier's breakdown and the run's
provenance, and replaying a log regenerates each problem from its seed and re-scores the completion exactly.
The result files that the tables and the figure read sit beside their logs.
The base models are public checkpoints, loaded from pinned snapshots (\pv{model.base} and \pv{model.llama}).
The trained LoRA adapters are not distributed with this version: the repository holds each adapter's configuration
and the fingerprint recorded at the end of its training, not its weights.
Re-running a model's evaluation therefore means retraining the adapter from its configuration and seed, and GPU
training is not bit-for-bit deterministic, so a retrained adapter can differ from the one evaluated here.

% D113 sec. 8 (manifests, FROZEN.md, the loader guard); D67 and D74 (the noise-free sets and X2's candidates); D147
% (the J set); generated table tab:frozen.
\paragraph{Frozen evaluation sets.}
Table~\ref{tab:frozen} lists every frozen evaluation set with the manifest that holds its SHA-256 checksum.

% docs/hypothesis_ledger.yaml (D185, written from the FINDINGS entries it cites); PREREGISTRATION.md; D141 (family H
% not run); docs/decision_timeline.md (amendments after a result: D47, D54, D57, D92, D135).
\paragraph{The pre-registration ledger.}
Table~\ref{tab:ledger} marks every pre-registered hypothesis or test family with its verdict (supported or not,
rejected or not rejected, confirmed, or not run), the entry that registered it and the finding that reports it.
Amendments made after a result are in Table~\ref{tab:deviations} and Appendix~\ref{app:timeline}.

% GENERATED by scripts/paper/tables.py; do not edit. Rebuild: uv run python scripts/paper/build.py
% artifact: table:frozen
% sources (path @ git blob):
%   data/eval_sets/1.1.0+c735a0d4/manifest.json @ 16d096237235
%   data/eval_sets/1.1.0+c735a0d4/x2_candidates_manifest.json @ 79f42fb18370
%   data/eval_sets/2.0.0+bdbb2279/confirmatory/a6_pairs_manifest.json @ a4667a2fef54
%   data/eval_sets/2.0.0+bdbb2279/confirmatory/a6b_pairs_manifest.json @ f42552284794
%   data/eval_sets/2.0.0+bdbb2279/confirmatory/a6c_pairs_manifest.json @ e413b9b471bf
%   data/eval_sets/2.0.0+bdbb2279/confirmatory/r5_6/a6b_pairs_manifest.json @ 33486465e0b3
%   data/eval_sets/2.0.0+bdbb2279/confirmatory/r5_6/a6c_pairs_manifest.json @ 3e36aa054364
\begin{table}[t]
\centering
\footnotesize
\caption{The frozen evaluation sets.}
\label{tab:frozen}
\begin{tabular}{@{}l l r l l@{}}
\toprule
set & generator & problems & used in & SHA-256 \\
\midrule
in-distribution, noise-free & 1.1.0+c735a0d4 & \num{200} & Results I & \texttt{970d71a7761cb088\ldots{}} \\ \multicolumn{5}{@{}l@{}}{\quad\scriptsize\path{data/eval_sets/1.1.0+c735a0d4/manifest.json}} \\
depth 5, noise-free & 1.1.0+c735a0d4 & \num{200} & Results I & \texttt{0e7f34ed11437133\ldots{}} \\ \multicolumn{5}{@{}l@{}}{\quad\scriptsize\path{data/eval_sets/1.1.0+c735a0d4/manifest.json}} \\
X2 candidates & 1.1.0+c735a0d4 & \num{1200} & Results II & \texttt{cd703f0d15cd69e4\ldots{}} \\ \multicolumn{5}{@{}l@{}}{\quad\scriptsize\path{data/eval_sets/1.1.0+c735a0d4/x2_candidates_manifest.json}} \\
R4-5 A6 & 2.0.0+bdbb2279 & \num{1020} & Results III & \texttt{e803a5676c66fc04\ldots{}} \\ \multicolumn{5}{@{}l@{}}{\quad\scriptsize\path{data/eval_sets/2.0.0+bdbb2279/confirmatory/a6_pairs_manifest.json}} \\
R4-5 A6b & 2.0.0+bdbb2279 & \num{1020} & Results III & \texttt{a4f6d81f73dd744e\ldots{}} \\ \multicolumn{5}{@{}l@{}}{\quad\scriptsize\path{data/eval_sets/2.0.0+bdbb2279/confirmatory/a6b_pairs_manifest.json}} \\
R4-5 A6c & 2.0.0+bdbb2279 & \num{1020} & Results III & \texttt{95ba42801c89b11c\ldots{}} \\ \multicolumn{5}{@{}l@{}}{\quad\scriptsize\path{data/eval_sets/2.0.0+bdbb2279/confirmatory/a6c_pairs_manifest.json}} \\
J sensitivity & 2.0.0+bdbb2279 & \num{816} & Generalization & \texttt{42f2b3d1ffc7e8d1\ldots{}} \\ \multicolumn{5}{@{}l@{}}{\quad\scriptsize\path{data/eval_sets/2.0.0+bdbb2279/confirmatory/r5_6/a6c_pairs_manifest.json}} \\
J invariance & 2.0.0+bdbb2279 & \num{816} & Generalization & \texttt{395f04dc06d53293\ldots{}} \\ \multicolumn{5}{@{}l@{}}{\quad\scriptsize\path{data/eval_sets/2.0.0+bdbb2279/confirmatory/r5_6/a6b_pairs_manifest.json}} \\
\bottomrule
\end{tabular}
\par\smallskip\begin{minipage}{0.95\linewidth}\footnotesize Each set is regenerated from its seeds and compared with its file before any generation; the loader refuses the confirmatory sets without an explicit flag (D113). The first 16 hexadecimal digits of each file's SHA-256 are shown; the full digest is in the manifest named under it.\end{minipage}
\end{table}

% GENERATED by scripts/paper/tables.py; do not edit. Rebuild: uv run python scripts/paper/build.py
% artifact: table:ledger
% sources (path @ git blob):
%   docs/hypothesis_ledger.yaml @ 31fe9224941f
\begingroup\footnotesize
\begin{longtable}{@{}L{1.7cm} L{1.8cm} L{3.7cm} L{1.9cm} L{3.7cm} L{1.2cm}@{}}
\caption{The pre-registration ledger.}\label{tab:ledger} \\
\toprule
tests & entry & question & label & verdict & finding \\
\midrule
\endfirsthead
\toprule
tests & entry & question & label & verdict & finding \\
\midrule
\endhead
\midrule
\multicolumn{6}{@{}p{15.4cm}@{}}{Every pre-registered hypothesis or test family, in the order of the pre-registration and the decision log, with its verdict as the cited FINDINGS entry states it. Supported or rejected: the registered one-sided test passed (Holm-adjusted where the family is); not rejected: it did not, which is not evidence of equivalence. Source: docs/hypothesis\_ledger.yaml.} \\
\endlastfoot
P1--P3 & pre-registration, D45 & the warm-up model's hidden state predicts its own error better than its stated $\sigma$ & confirmatory & uninterpretable; P1--P3 pass Holm, but the registered random-init sanity check fails & \#7 \\
S1a, S1b & pre-registration, D62 & split credit raises $1\sigma$ coverage above the warm-up model (S1a) and uniform credit (S1b) & confirmatory & supported & \#16, \#18 \\
S2 & pre-registration, D62 & split credit's $\sigma$ ranks errors better than uniform credit's & confirmatory & not supported (an earlier narrow pass came from the bf16 merge and is withdrawn) & \#13, \#16 \\
S3 & pre-registration, D62 & split credit keeps accuracy & confirmatory & supported & \#16 \\
N1--N3 & D50 (A5) & split credit's $\sigma$ beats the text baseline A5 on the same answers & confirmatory & not supported on any seed & \#16 \\
gate (A3) & D50, D54, D57 & the generator's $\sigma^\star$ intervals cover the clean truth at nominal rates & confirmatory, amended twice after failing & passed on the final run & \#9, \#11, \#12 \\
E1a-P, E1b-P & D52 & with the answer fixed, RL on $\sigma$ alone beats the warm-up model's $\sigma$ & confirmatory (pilot) & supported & \#14, \#16 \\
E1a-N, E1b-N & D52 & the same $\sigma$ beats the A5 text baseline & confirmatory (pilot) & not supported & \#14, \#16 \\
X1 & D73 & a shallow text model of $\sigma^\star$ propagates on A6 & confirmatory & labelled partial propagation on A6, slope far below one; does not transfer to A6c & \#17, \#23 \\
X2 & D74 & a single-completion readout ranks a right answer above a wrong one to the same problem & confirmatory & stated $\sigma$ and text models not supported; the probe is supported but not distinguishable from its random-init control & \#19 \\
X3 & D79 & the RL setup has the capacity to learn a $\sigma$ that ranks errors (distillation) & confirmatory (pilot) & capacity likely, between problems only; at chance within problems & \#20 \\
X4 & D89, D91 & the model's own verdict on its answer ranks right above wrong within problems & confirmatory & not supported & \#21 \\
E1c, E1d & D94, D101 & batch-std scaling (E1c) or a larger learning rate (E1d) lets $\sigma$-only RL rank errors & confirmatory (pilot); reported as exploratory (R5-7) & E1d supported (between problems); E1c not supported & \#25 \\
T1--T3 & D102 & E1d's $\sigma$ ranks right above wrong within problems & confirmatory & not supported & \#26 \\
B1--B5 & D104 & supervised $\sigma^\star$ (stage B) propagates and beats RSS and X1 & confirmatory & not supported & \#27 \\
C1--C6 & D106 & $\sigma$-only RL with the predictive reward (stage C) propagates and beats its start & confirmatory (pilot) & C6 (beats its start) supported; C1--C5 not supported & \#28 \\
D1--D7 & D109 & an explicit propagation line (stage D) propagates and beats the references & confirmatory (pilot) & D6, D7 supported; D1--D5 not supported & \#29 \\
R4-3 go rule & D113 & D-pos-v2 meets the propagation criterion on development data (go / no-go) & confirmatory (pilot) & GO on all three $\sigma$ values; no score gain over RSS & \#30 \\
MC-A stopping rule & D113 & re-asking the model on perturbed inputs beats RSS & diagnostic & not supported; the stopping rule fires & \#31 \\
P (K1, C2, C3) & D118, D125 & primary: D-pos-v2's answer $\sigma$ propagates on the frozen set & confirmatory (primary) & confirmed & \#32 \\
K1--K5 & D118 & the three D-pos-v2 $\sigma$s and B-v2 propagate; D-pos-v2 scores above B-v2 & confirmatory & all rejected under Holm (supported) & \#32 \\
K6--K8 & D118 & D-pos-v2 beats RSS (K6) and X1 (K7); B-v2 beats RSS (K8) & confirmatory & not rejected & \#32 \\
M1--M5 & D119, D122 & the model's own program (MC-B) beats RSS, X1, D-pos-v2 and MC-A, and propagates & confirmatory & not rejected & \#33 \\
cross-model go rule & D137, D139, D146 & the R4-3 criterion on a second model family at one seed & confirmatory (pilot) & NO-GO; the sensitivity criterion passed, invariance could not be measured & \#38 \\
P$_J$ (J1, C2, C3) & D147 & D-pos-v2's $\sigma$ propagates through Jacobians outside the training range of trained chains & confirmatory & confirmed (a clamped in-range propagator reaches most of the contrast) & \#39 \\
J1, J2 & D147 & the J-sensitivity slope is positive; explicit exceeds implicit & confirmatory & rejected under Holm (supported) & \#39 \\
J3--J6 & D147 & either arm beats RSS or X1 in expected score on the J set & confirmatory & not rejected & \#39 \\
H1--H5 & pre-registration, D141 & the stated-$\sigma$ model with a depth and noise curriculum (vol beta, conditional calibration, sources, decision value, stress) & not run & not run; the model they were written for was never built & -- \\
\end{longtable}
\endgroup


\section{The bf16 LoRA-merge pitfall}
\label{app:bf16}

% FINDINGS #15 (D65; exploratory); results/precision_check/merge_rounding.json (weights only); FINDINGS #4's R2-2
% note and D70, D71 (the SFT merge: every RL run started from the merged model; SFT-only evaluated as the runs saw it,
% #16).
\paragraph{What happened.}
Reinforcement-learning adapters were trained on the \pv{model.bits}-bit base with the LoRA unmerged, but an early
evaluation path merged each adapter into a bf16 copy of the warm-up model before generating.
Merging rounds every updated weight to the nearest bf16 value.
The RL updates were far smaller than one bf16 step of the base weights (per adapter, a median of \pv{bf.ulp} of a
step), so on the \pv{bf.nadapters} RL adapters checked, \pv{bf.erased} of each update's entries were rounded to
exactly zero.
The warm-up adapter, whose updates are larger, lost \pv{bf.sft.erased} when it was merged; every reinforcement-learning
run started from that merged model, and the warm-up model is evaluated as the runs saw it (merged, then quantized), so
this loss defines the warm-up model rather than distorting comparisons between runs.

% FINDINGS #15 (the training-vs-evaluation sigma gap; D65's rule, stated before the unmerged evaluations);
% results/precision_check/report.json (split credit, the first seed at the matched step).
\paragraph{How it was caught.}
The evaluated $\sigma$ was several times narrower than the $\sigma$ the same policy stated during training.
For the first split-credit seed, the median relative $\sigma$ was \pv{bf.sigma.train} in its last training steps and
\pv{bf.sigma.merged} when sampled from the merged model.
A check stated before its data (D65) evaluated the adapters as trained, unmerged on the quantized base.
Sampled, the same adapter then stated \pv{bf.sigma.unmerged}, in line with training, and its greedy $1\sigma$
coverage rose from \pv{bf.cov.merged} to \pv{bf.cov.unmerged}.
Only \pv{bf.identical} of the paired greedy completions were identical between the two ways of applying the adapter,
and the median per-problem ratio of the stated $\sigma$ was \pv{bf.ratio}.
Weight-level diagnostics then placed the cause in the merge rather than in the quantized base.

% D67, D68 (the corrected, unmerged evaluations become primary; merged results kept, labelled superseded); FINDINGS
% #13, #14 (superseded), #16 (S2's verdict changed); docs/paper_outline.md claims to avoid (scope of the bf16 claim).
\paragraph{What was done.}
From then on every adapter was evaluated as trained, unmerged on the same quantized base, and the corrected
evaluations became the primary results for the families they affected.
The merged results stay in the record, labelled superseded, including one verdict that changed (S2,
Table~\ref{tab:ledger}).
Every result in this paper uses the corrected evaluations.
The lesson is narrow: merging these small LoRA updates into the bf16 warm-up model erased most of them in our
pipeline; we make no claim about bf16 merging in general.

\section{The decision timeline}
\label{app:timeline}

% docs/decision_timeline.md (R5-8); D185 (R7-3 C); the deviations table (Section sec:deviations); D47, D54, D57, D92,
% D135 (the amendments made after a result).
Table~\ref{tab:timeline} lists every pre-registration and amendment in date order, with its first commit and its
analysis label: confirmatory (committed before its data and run as registered), a confirmatory pilot (the same at one
seed, which decided what ran next), or exploratory (post hoc, descriptive or diagnostic).
Results from before the pre-registration are exploratory.
Every amendment or deviation that touched a registered analysis is also in Table~\ref{tab:deviations}
(Section~\ref{sec:deviations}), with its timing relative to the data it affects.
Those made after a result are the generator gate's two amendments, the exploratory baselines of family P, the X2
robustness checks and the bounds of the RSS fallback.

% GENERATED by scripts/paper/tables.py; do not edit. Rebuild: uv run python scripts/paper/build.py
% artifact: table:timeline
% sources (path @ git blob):
%   docs/decision_timeline.md @ 8180e2e3c3ca
\begingroup\footnotesize
\begin{longtable}{@{}L{1.8cm} l L{5.0cm} L{3.1cm} L{3.1cm}@{}}
\caption{The decision timeline.}\label{tab:timeline} \\
\toprule
entry & date & what it fixed & kind & analysis \\
\midrule
\endfirsthead
\toprule
entry & date & what it fixed & kind & analysis \\
\midrule
\endhead
\midrule
\multicolumn{5}{@{}p{15.4cm}@{}}{Every pre-registration and amendment in date order (IST, 2026), with its first commit, what it fixed, its timing relative to the data it governs, and the analysis label. Confirmatory: committed before its data and run as registered; pilot: the same at one seed; exploratory: post hoc, descriptive or diagnostic. Results before the pre-registration (FINDINGS \#1--\#6) are exploratory. Source: docs/decision\_timeline.md.} \\
\endlastfoot
pre-registration \texttt{f601bb7} & 10-14 09:21 & Families P (does the SFT model know its error), S (split credit), H (H1--H5, phases 10--14); Holm within each family & pre-registration & P, S: \textbf{confirmatory}. H: \textbf{not run} (see the end) \\
D45, D46 \texttt{b662ab6} & 10-14 10:41 & Family P's probe implementation; family S's split-credit trainer & implementation, before results & confirmatory (P, S) \\
D47 \texttt{b662ab6} & 10-14 10:41 & P's pre-registered sanity check failed; exploratory baselines added afterwards & amendment, \textbf{after} P's result & exploratory (post-hoc) \\
D48 \texttt{f58eb20} & 10-14 10:42 & Scope of the rounding audit (Phase 10 step 2) & analysis plan & exploratory (descriptive) \\
D49 \texttt{6eeef4b} & 10-14 11:54 & Two post-hoc calibration baselines; RL calibration judged against the text-feature one & amendment, before RL calibration results & comparator rule (superseded by D50/A5) \\
D50 \texttt{5f2a846} & 10-14 11:59 & Review round 1, A1--A8: two estimands; the $\sigma$$^{*}$ definition; the generator self-consistency gate (A3); statistics; nested A5 baselines and family N (A5); A6 counterfactual pairs; cuts to family H (A7); the headline question & amendment, before any result it affects & gate and family N: confirmatory \\
D52 \texttt{fbb29a7} & 10-14 12:08 & E1, the fixed-$\mu$ $\sigma$-only GRPO pilot & pre-registration & confirmatory (pilot) \\
D54 \texttt{8acdba3} & 10-14 12:43 & Noise capped at 5\%; 10\% becomes a heavy-tail regime & amendment, \textbf{after} the gate failed (\#9) & gate rerun \\
D57 \texttt{5202b1f} & 10-14 12:56 & A3's cell-share rule replaced by a miss-count test; the next rerun is FINAL & amendment, \textbf{after} \#11, before the final rerun & gate (confirmatory, with two disclosed amendments) \\
D58 \texttt{113eb72} & 10-14 13:00 & Construction of the A6 sensitivity pairs & analysis plan, before any model on A6 & the A6 slope tests \\
D59 \texttt{0225b0a} & 10-14 13:11 & A6b invariance pairs (the slope must be $\equiv$ 0) & pre-registration & confirmatory (C3 of R4-3 uses it) \\
D61 \texttt{1203fb4} & 10-16 14:34 & A6 slope labels ("partial / full propagation") & amendment, before any A6 model result & labels \\
D62 \texttt{7a8f9f1} & 10-16 19:12 & Split-credit seeds 1 and 2; the seed-mean S1--S3 test; greedy vs sampled $\sigma$ & amendment to family S, before seed-1/2 data & S: confirmatory; greedy vs sampled: exploratory \\
D64 \texttt{9578ee6} & 10-17 21:57 & The label "anti-propagation" for a CI below 0 & amendment, before any A6 model result & labels \\
D65 \texttt{deff739} & 10-17 22:07 & RL adapters evaluated on the 4-bit training base & analysis plan, before any 4-bit eval & exploratory \\
D67 \texttt{bb7ff3c} & 10-18 00:06 & The bf16 merge was a bug: corrected (unmerged, 4-bit) evals become PRIMARY for families S, N, E; merged results kept, labelled superseded & amendment, before any corrected result & confirmatory (S, N, E re-run) \\
D69 \texttt{92b868d} & 10-18 09:48 & Post-hoc baselines re-fitted on 4-bit SFT calibration completions & amendment, before any such completion & comparator \\
D71 \texttt{55c8947} & 10-18 09:50 & Review round 2, R2-1--R2-7: evaluation hygiene and the reload fingerprint (R2-1); SFT revalidation (R2-2); what $\sigma$ should and should not predict (R2-3); seed-level statistics (R2-4); a publication criterion (R2-5); uniform-credit seeds 1 and 2 (R2-6); citations to verify before use (R2-7) & amendment, before any result it affects & R2-6: exploratory (robustness) \\
D72 \texttt{3f70b9a} & 10-18 10:06 & R2-1's fingerprint and reload-check tolerances & analysis plan, before any fingerprint & reproducibility check \\
D73 \texttt{3eaab53} & 10-18 10:07 & X1, a shallow text model of $\sigma$$^{*}$ on A6 / A6b & pre-registration & confirmatory \\
D74 \texttt{d9a36c9} & 10-18 22:12 & X2: within-problem correctness pairs & pre-registration & confirmatory \\
D77 \texttt{42bd61e} & 10-19 10:31 & R2-1: below the log-prob floor, the $\sigma$-token criterion decides; the check is enforced & amendment (future fingerprints only) & reproducibility check \\
D79 \texttt{de83afd} & 10-19 10:40 & X3, the distillation capacity test & pre-registration & confirmatory (pilot) \\
D89 \texttt{eddda39} & 10-20 08:35 & X4, self-verification on the X2 pairs & pre-registration & confirmatory \\
D91 \texttt{4f016c6} & 10-20 23:22 & X4: a direct mode, a perturbed control, the cluster bootstrap, token cost & amendment, before any X4 data & confirmatory \\
D92 \texttt{5b868f3} & 10-20 23:38 & X2 robustness checks, defined before computing & analysis plan, \textbf{after} X2's result & exploratory (post-hoc) \\
D93 \texttt{1faf388} & 10-20 23:45 & Review round 3, R3-1--R3-10: X4 additions; X2 robustness; optimization pilots; Phase 12 scope = aleatoric propagation (R3-4); A6c (R3-5); tolerances; statistics; claim strength & amendment, before any result it affects & -- \\
D94 \texttt{43cc2af} & 10-20 23:49 & E1c and E1d, the R3-3 optimization pilots & pre-registration & confirmatory (pilot). The paper states E1d as exploratory (R5-7) \\
D95 \texttt{4ff29f2} & 10-20 23:53 & A6c noise-allocation pairs and the noise heuristics: construction and analysis & analysis plan, before the set is built & exploratory (descriptive bars) \\
D99 \texttt{7792502} & 10-21 06:07 & X1 on A6c & analysis plan, before any prediction & exploratory (descriptive) \\
D100 \texttt{7792502} & 10-21 06:07 & Stage A: the scoring audit and the oracle-centred diagnostic & analysis plan, before computing & exploratory (descriptive) \\
D101 \texttt{b24a782} & 10-21 10:44 & D94's decision rule applied: E1d's learning rate is adopted & decision under a registered rule & -- \\
D102 \texttt{2970d07} & 10-21 23:14 & E1d's within-problem test on the X2 pairs (T1--T3) & pre-registration & confirmatory \\
D104 \texttt{dcd9b12} & 10-21 23:20 & Stage B, the supervised $\sigma$$^{*}$ ceiling, 3 seeds & pre-registration & confirmatory \\
D106 \texttt{34738a1} & 10-21 23:36 & Stage C pilot: $\sigma$-only RL with verifier 1.3.0's $\sigma$ reward & pre-registration & confirmatory (pilot) \\
D109 \texttt{ff59b4a} & 10-22 23:44 & Stage D: explicit propagation, a positive control and a format-only ablation, 1 seed per arm & pre-registration & confirmatory (pilot) \\
D112 \texttt{87a5933} & 10-25 17:13 & Review round 4, R4-1--R4-9: D-pos-v2 vs B-v2 on mixed-level data, MC-A, the frozen set, framing (R4-8) & amendment, before any result it affects & -- \\
D113 \texttt{e504498} & 10-25 17:20 & R4-1--R4-5: D-pos-v2 and B-v2, the three-$\sigma$ grading, the stop/go rule (K1 $\wedge$ C2 $\wedge$ C3), MC-A, the frozen confirmatory set & pre-registration & confirmatory (pilot: the go rule); MC-A a registered diagnostic \\
D114 \texttt{a4b0a2a} & 10-25 17:31 & Two deviations found while building: eval budget 1,280 tokens; 6 rows dropped & deviation, before any R4 output & -- \\
D116 \texttt{dc54969} & 10-25 19:07 & D-pos-v2's eval OOM'd in its backoff; fixed and rerun at batch 16 & deviation, before any completion was scored & -- \\
D118 \texttt{f64443e} & 10-26 07:26 & \textbf{The R4-3 confirmatory run:} D-pos-v2 and B-v2, 3 seeds each, on the frozen R4-5 set; K1--K8, Holm & pre-registration & \textbf{confirmatory (primary)} \\
D119 \texttt{d3de520} & 10-26 07:28 & MC-B: the model writes the problem's program; a sandbox runs it & pre-registration & confirmatory \\
D122 \texttt{f43ccb8} & 10-26 07:52 & MC-B: two deviations found while building & deviation, before any output & -- \\
D123 \texttt{e93ab45} & 10-26 08:00 & D118's analysis code built and checked; X1's confirmatory predictions computed & analysis frozen before data & confirmatory \\
D125 \texttt{5ecdeea} & 10-26 08:34 & D118 \S{}6 amended: the primary P is D113's full criterion (K1 $\wedge$ C2 $\wedge$ C3) & amendment, before any confirmatory result & \textbf{confirmatory (primary)} \\
D131 \texttt{c9af378} & 10-28 02:52 & D-pos-v2 seed 1's eval hung after a GPU driver reset; rerun on identical code & deviation (operational), before the analysis & -- \\
D133 \texttt{be0d338} & 10-28 10:04 & Review round 5, R5-1--R5-10: framing; tail autopsy (R5-2); hybrid oracle (R5-3); RSS fallback (R5-4); cross-model replication (R5-5); cuts (R5-7); reporting rules (R5-8) & amendment, \textbf{after} \#32 / \#33, before anything it asks for & R5-2, R5-8: exploratory (descriptive). R5-3: exploratory (diagnostic) \\
D135 \texttt{bad601f} & 10-28 10:24 & R5-4's RSS-fallback bounds, $\sigma$/RSS $\in$ [0.333, 5], chosen on the development sets and frozen before one application to the confirmatory logs & analysis plan; grid designed \textbf{after} seeing \#35 & exploratory (post-hoc) \\
D137 \texttt{7892e89} & 10-28 18:07 & The cross-model pilot (R5-5): D-pos-v2 and B-v2 on Llama-3.2-3B-Instruct, seed 0, frozen A6c + A6b; go rule = R4-3's criterion on D-pos-v2; 3 seeds only on GO & pre-registration & confirmatory (pilot) \\
D138 \texttt{b50cae1} & 10-28 18:46 & R2-1's fingerprint for the Llama base records at sampling batch 4 & deviation (implementation), before any pilot output & -- \\
D139 \texttt{9296d7b} & 10-29 00:13 & D137's go rule runs as soon as the D-pos-v2 eval exists (same code, same logs); SFT seed 1 only on GO & amendment (timing), before the go rule runs & confirmatory (pilot) \\
D147 \texttt{8405f0f} & 10-30 14:39 & \textbf{R5-6 part J:} Jacobians outside the training range of trained chains; the frozen J set (sensitivity + dimensional-scaling invariance pairs); D-pos-v2 and B-v2 $\times$ 3 seeds; P\_J = J1 $\wedge$ C2 $\wedge$ C3, Holm over J1--J6 & pre-registration & \textbf{confirmatory} \\
D148 \texttt{8405f0f} & 10-30 14:39 & Experiment freeze: R5-6 part J is the paper's last experiment; everything else is future work & scope decision & -- \\
D149 \texttt{5f17510} & 10-30 15:00 & R5-6 J build record: support, constructions, frozen set (all cells full), X1, the analysis checked on development files & build, before any J completion & confirmatory \\
\end{longtable}
\endgroup

